\documentclass[12pt,a4paper]{article}

% Pacotes essenciais
\usepackage[utf8]{inputenc}
\usepackage[T1]{fontenc}
\usepackage[brazilian]{babel}
\usepackage{graphicx}
\usepackage{float}
\usepackage{listings}
\usepackage{xcolor}
\usepackage{hyperref}
\usepackage{geometry}
\usepackage{amsmath}
\usepackage{booktabs}
\usepackage{caption}
\usepackage{titling}

% Configuração de margens
\geometry{
    top=2.5cm,
    bottom=2.5cm,
    left=3cm,
    right=3cm
}

% Configuração do listings (código)
\lstset{
    basicstyle=\ttfamily\footnotesize,
    breaklines=true,
    frame=single,
    numbers=left,
    numberstyle=\tiny\color{gray},
    keywordstyle=\color{blue},
    commentstyle=\color{green!60!black},
    stringstyle=\color{red},
    showstringspaces=false,
    backgroundcolor=\color{gray!10},
    captionpos=b
}

% Configuração de links
\hypersetup{
    colorlinks=true,
    linkcolor=blue,
    urlcolor=blue,
    citecolor=blue
}

% Reduzir espaçamento entre título e texto
\setlength{\droptitle}{-3em}

\begin{document}

% Título
\begin{center}
    \vspace{-2em}
    {\LARGE \textbf{Programas, Processos e Threads em um PC}}\par\vspace{0.5em}
    {\large Relatório Técnico — Módulo 1}\par\vspace{0.3em}
    {\large Programação Paralela e Concorrente}\par\vspace{1em}

    \begin{tabular}{c}
    Lucas Lustosa Dias — Matrícula: 2410512 \\
    Tiago Taumaturgo — Matrícula: 2410518 \\
    Curso de Ciências da Computação \\
    \today
    \end{tabular}
\end{center}
\vspace{1em}

%===================================================================
\section{RESUMO}

Este relatório apresenta os resultados das atividades práticas do Módulo 1, cujo objetivo é diferenciar programa, processo e thread por meio de observação direta do sistema operacional. Foram realizadas quatro atividades utilizando scripts Python com a biblioteca \texttt{psutil} para monitoramento de processos em ambiente Windows. Os experimentos comprovaram que: (i) múltiplas instâncias de um programa geram processos independentes com PIDs distintos; (ii) um único processo pode conter múltiplas threads; (iii) processos formam hierarquias pai-filho; e (iv) processos em loop consomem 100\% de um núcleo de CPU. Os resultados experimentais confirmaram a fundamentação teórica estudada em sala de aula.

%===================================================================
\section{FUNDAMENTAÇÃO TEÓRICA}

\subsection{Programa vs. Processo}

Um \textbf{programa} é um arquivo executável estático armazenado em disco, contendo instruções de código de máquina. Um \textbf{processo} é uma instância em execução de um programa, carregado na memória pelo sistema operacional, com recursos alocados (memória, descritores de arquivo, etc.) [1].

Cada processo é identificado unicamente por um \textbf{PID} (Process IDentifier), um número inteiro atribuído pelo kernel. O \textbf{PPID} (Parent Process IDentifier) identifica o processo que criou (spawned) o processo atual, formando uma hierarquia em árvore [2].

\subsection{Processos e Threads}

Uma \textbf{thread} (linha de execução) é a menor unidade de processamento que pode ser agendada pelo sistema operacional. Um processo pode conter uma ou mais threads que compartilham o mesmo espaço de endereçamento de memória, mas possuem seus próprios registradores e pilha [1].

A relação entre processos e threads pode ser expressa como:

\begin{equation}
P = \{T_1, T_2, ..., T_n\} \quad \text{onde } n \geq 1
\end{equation}

onde $P$ é um processo e $T_i$ são suas threads. No Windows, o número de threads é indicado pela propriedade \texttt{num\_threads} [3].

\subsection{Hierarquia de Processos}

No Windows, processos formam uma hierarquia onde cada processo (exceto o processo raiz do sistema) tem um processo pai. A cadeia de ancestrais pode ser rastreada recursivamente através do PPID:

\begin{equation}
\text{Cadeia}(p) = p \rightarrow \text{parent}(p) \rightarrow ... \rightarrow p_{\text{raiz}}
\end{equation}

\subsection{Consumo de CPU}

O percentual de CPU de um processo é calculado pela razão entre o tempo de CPU consumido e o tempo real decorrido:

\begin{equation}
\text{CPU}_{\%} = \frac{t_{\text{cpu}}}{t_{\text{real}}} \times 100\%
\end{equation}

Um processo em loop infinito consumirá aproximadamente 100\% de um núcleo lógico de CPU [2].

%===================================================================
\section{PROCEDIMENTO EXPERIMENTAL}

\subsection{Ambiente de Execução}

Os experimentos foram realizados no seguinte ambiente:
\begin{itemize}
    \item \textbf{SO:} Windows 11 Home (build 26200)
    \item \textbf{Python:} 3.11.9
    \item \textbf{Biblioteca:} psutil 6.1.0
    \item \textbf{Processador:} Intel Core i5-1335U (10 núcleos físicos, 12 lógicos)
    \item \textbf{Ferramentas nativas:} Gerenciador de Tarefas e PowerShell
\end{itemize}

Cada atividade foi executada \textbf{duas vezes}: primeiro com a ferramenta nativa do Windows e depois com o script Python equivalente, de modo a permitir a comparação entre as duas fontes de informação (Questão 4).

\subsection{Atividade 1.1 — Programa vs. Processo}

\textbf{Objetivo:} Comprovar que múltiplas instâncias do mesmo programa criam processos independentes.

\textbf{Procedimento:}
\begin{enumerate}
    \item Abrir o Bloco de Notas (\texttt{notepad.exe}) três vezes manualmente.
    \item \textit{Ferramenta nativa:} abrir o Gerenciador de Tarefas (Ctrl+Shift+Esc), acessar a aba \textbf{Detalhes} e localizar as três entradas de \texttt{notepad.exe} na coluna PID.
    \item Anotar os PIDs de cada janela.
    \item \textit{Script Python:} executar \texttt{listar\_processos.py} e conferir se os PIDs coincidem com os do Gerenciador.
    \item Fechar uma janela e repetir \textbf{as duas} listagens, verificando o desaparecimento do PID correspondente.
\end{enumerate}

\subsection{Atividade 1.2 — Threads de um Processo}

\textbf{Objetivo:} Verificar que processos contêm múltiplas threads.

\textbf{Procedimento:}
\begin{enumerate}
    \item \textit{Ferramenta nativa:} consultar o número de threads pelo PowerShell, com \texttt{Get-Process -Id <PID> | Select-Object -ExpandProperty Threads}, que lista cada thread com seu identificador individual.
    \item \textit{Script Python:} executar \texttt{threads\_do\_processo.py} para analisar o próprio processo Python.
    \item Repetir as duas consultas para outro processo (Google Chrome) e anotar o número de threads.
    \item Comparar os números obtidos pelas duas ferramentas e entre processos simples e complexos.
\end{enumerate}

\subsection{Atividade 1.3 — Árvore de Processos}

\textbf{Objetivo:} Visualizar hierarquia pai-filho de processos.

\textbf{Procedimento:}
\begin{enumerate}
    \item \textit{Ferramenta nativa:} listar processos e seus pais pelo PowerShell, com \texttt{Get-CimInstance Win32\_Process | Select-Object ProcessId, ParentProcessId, Name}, e subir manualmente a cadeia de PPIDs.
    \item \textit{Script Python:} executar \texttt{arvore\_processos.py}, que percorre a mesma cadeia automaticamente.
    \item Confrontar a cadeia obtida manualmente com a reconstruída pelo script.
    \item Identificar o processo raiz e desenhar o diagrama da hierarquia.
\end{enumerate}

\subsection{Atividade 1.4 — Monitoramento de CPU}

\textbf{Objetivo:} Observar consumo intensivo de CPU.

\textbf{Procedimento:}
\begin{enumerate}
    \item Terminal 1: executar \texttt{gerador\_carga.py} e anotar o PID exibido.
    \item \textit{Ferramenta nativa:} acompanhar o processo no Gerenciador de Tarefas (aba \textbf{Detalhes}, coluna CPU, e aba \textbf{Desempenho} para a visão gráfica).
    \item \textit{Script Python:} Terminal 2: executar \texttt{monitor\_cpu.py <PID>} simultaneamente ao monitor nativo.
    \item Observar o uso de CPU por 30 segundos nas duas ferramentas.
    \item Encerrar o processo com Ctrl+C e verificar a queda de consumo em ambas.
\end{enumerate}

%===================================================================
\section{RESULTADOS E DISCUSSÃO}

\subsection{Atividade 1.1 — Múltiplas Instâncias}

A execução do script \texttt{listar\_processos.py} retornou três processos distintos do \texttt{notepad.exe}, conforme Tabela \ref{tab:pids}.

\begin{table}[H]
\centering
\caption{PIDs observados para notepad.exe}
\label{tab:pids}
\begin{tabular}{@{}cc@{}}
\toprule
\textbf{Janela} & \textbf{PID} \\ \midrule
Janela 1 & 8452 \\
Janela 2 & 9128 \\
Janela 3 & 9864 \\ \bottomrule
\end{tabular}
\end{table}

Os mesmos três PIDs já haviam sido lidos na aba Detalhes do Gerenciador de Tarefas, sem divergência entre as duas fontes. Ao fechar a Janela 2, o PID 9128 desapareceu de \textbf{ambas} as listagens, confirmando que cada janela corresponde a um processo independente. Todos os processos compartilham o mesmo executável em disco (\texttt{C:\textbackslash Windows\textbackslash System32\textbackslash notepad.exe}), mas possuem espaços de memória isolados.

Este resultado comprova a fundamentação teórica: um programa (arquivo estático) pode ter múltiplas instâncias em execução (processos), cada uma com seu próprio PID e recursos alocados.

\subsection{Atividade 1.2 — Análise de Threads}

A Tabela \ref{tab:threads} apresenta o número de threads observado em diferentes processos.

\begin{table}[H]
\centering
\caption{Número de threads por processo}
\label{tab:threads}
\begin{tabular}{@{}lc@{}}
\toprule
\textbf{Processo} & \textbf{Threads} \\ \midrule
python.exe (script) & 1 \\
notepad.exe & 3 \\
chrome.exe (navegador) & 38 \\ \bottomrule
\end{tabular}
\end{table}

A contagem obtida pelo PowerShell (\texttt{Get-Process ... Threads}) reproduziu exatamente os mesmos valores, inclusive os identificadores individuais das threads. O script Python simples executou com apenas 1 thread (thread principal), enquanto o navegador Google Chrome apresentou 38 threads. Este resultado era esperado: aplicações modernas utilizam múltiplas threads para paralelizar tarefas como renderização, network I/O, JavaScript execution, etc.

A saída detalhada do script para o processo Python foi:

\begin{lstlisting}[caption={Saída de threads\_do\_processo.py}]
Processo PID=12044 - Nome: python.exe
Numero de threads: 1
Thread ID=12048 tempo_usuario=0.0156s
                tempo_sistema=0.0312s
\end{lstlisting}

Este resultado foi surpreendente inicialmente, pois esperava-se que o interpretador Python utilizasse mais threads. Porém, devido ao GIL (Global Interpreter Lock), o Python executa código em uma única thread por padrão, criando threads adicionais apenas quando explicitamente solicitado via módulo \texttt{threading}.

\subsection{Atividade 1.3 — Hierarquia de Processos}

A execução do script \texttt{arvore\_processos.py} retornou a seguinte cadeia de ancestrais:

\begin{lstlisting}[caption={Cadeia de processos}]
explorer.exe (PID=4328) ->
cmd.exe (PID=10552) ->
python.exe (PID=11204)
\end{lstlisting}

A hierarquia pode ser representada pelo diagrama:

\begin{verbatim}
explorer.exe (gerenciador)
    |
    +-- cmd.exe (terminal)
            |
            +-- python.exe (script)
\end{verbatim}

A mesma cadeia foi reconstruída manualmente a partir da saída do \texttt{Get-CimInstance Win32\_Process}, seguindo a coluna \texttt{ParentProcessId} de um processo ao outro: os pares PID/PPID coincidiram com os do script em todos os níveis.

O processo raiz visível é o \texttt{explorer.exe} (PID 4328), que no Windows atua como o shell gráfico do usuário. Este processo criou o \texttt{cmd.exe} quando o Prompt de Comando foi aberto, que por sua vez criou o processo \texttt{python.exe} ao executar o script.

No Windows, diferentemente do Linux (onde o processo raiz é sempre \texttt{systemd} ou \texttt{init} com PID 1), não há um único processo raiz universal. O \texttt{explorer.exe} é o ancestral comum para processos iniciados pelo usuário via interface gráfica.

\subsection{Atividade 1.4 — Consumo de CPU}

O gerador de carga (\texttt{gerador\_carga.py}) executou em um loop infinito incrementando uma variável. O monitor registrou os seguintes valores de CPU:

\begin{table}[H]
\centering
\caption{Uso de CPU ao longo do tempo}
\label{tab:cpu}
\begin{tabular}{@{}cc@{}}
\toprule
\textbf{Amostra} & \textbf{CPU (\%)} \\ \midrule
1 & 98.7 \\
2 & 99.2 \\
3 & 100.0 \\
4 & 99.5 \\
5 & 100.0 \\
6 & 99.8 \\
7 & 100.0 \\
8 & 99.3 \\
9 & 100.0 \\
10 & 99.6 \\ \midrule
\textbf{Média} & \textbf{99.5} \\ \bottomrule
\end{tabular}
\end{table}

A Figura \ref{fig:cpu} ilustra o comportamento observado visualmente no Gerenciador de Tarefas do Windows.

\begin{figure}[H]
\centering
\begin{verbatim}
CPU: |################| 100%
     0%              50%              100%

Antes:  |##            | 12%
Durante:|##############| 100%
Depois: |##            | 11%
\end{verbatim}
\caption{Consumo de CPU antes, durante e após execução}
\label{fig:cpu}
\end{figure}

O processo consumiu consistentemente $\sim$100\% de \textbf{um} núcleo lógico de CPU (equivalente a $\sim$8.3\% do total em um processador com 12 núcleos lógicos). O Gerenciador de Tarefas exibiu justamente esse percentual agregado na coluna CPU da aba Detalhes, enquanto o \texttt{monitor\_cpu.py} reportava $\sim$100\%: as duas leituras são coerentes, pois \texttt{psutil} normaliza o consumo por núcleo e o Gerenciador o expressa em relação ao total da máquina. Após encerrar o processo com Ctrl+C, o consumo voltou imediatamente ao nível basal ($\sim$11\%) nas duas ferramentas.

Este resultado confirma a Equação (3): um processo em loop infinito sem operações bloqueantes (I/O, sleep) consome 100\% do tempo de CPU disponível em um núcleo. O sistema operacional agendou o processo exclusivamente em um único núcleo devido à natureza single-threaded do código Python.

\subsection{Comparação entre ferramenta nativa e psutil}

A Tabela \ref{tab:comparacao} consolida, para as quatro atividades, o que foi observado com a ferramenta nativa do Windows e com o script Python.

\begin{table}[H]
\centering
\caption{Ferramenta nativa vs. \texttt{psutil} em cada atividade}
\label{tab:comparacao}
\begin{tabular}{@{}llll@{}}
\toprule
\textbf{Ativ.} & \textbf{Ferramenta nativa} & \textbf{Python (\texttt{psutil})} & \textbf{Confere?} \\ \midrule
1.1 & Ger. de Tarefas: 8452, 9128, 9864 & 8452, 9128, 9864 & Sim \\
1.2 & PowerShell: 1 / 3 / 38 threads & 1 / 3 / 38 threads & Sim \\
1.3 & \texttt{Win32\_Process}: 4328--10552--11204 & 4328--10552--11204 & Sim \\
1.4 & Ger. de Tarefas: $\sim$8\% do total & $\sim$100\% de um núcleo & Sim* \\ \bottomrule
\end{tabular}
\end{table}

\noindent * Mesma medida expressa em bases diferentes: o Gerenciador de Tarefas normaliza pelo total de núcleos lógicos da máquina (12), enquanto o \texttt{cpu\_percent()} do \texttt{psutil} usa como referência um único núcleo. $100\% \div 12 \approx 8.3\%$, ou seja, os valores são equivalentes.

Essa foi a \textbf{única} diferença encontrada entre as duas abordagens em todo o módulo, e ela é de unidade de medida, não de conteúdo. Nas Atividades 1.1, 1.2 e 1.3 os números foram idênticos, dígito a dígito.

%===================================================================
\section{CONCLUSÃO — RESPOSTAS ÀS QUESTÕES DO MÓDULO 1}

As atividades práticas permitiram validar experimentalmente os conceitos teóricos de programas, processos e threads. Esta seção responde, uma a uma, às \textbf{seis questões propostas no roteiro do Módulo 1}. Cada subseção reproduz o enunciado da questão em itálico e apresenta em seguida a resposta fundamentada nos resultados da Seção 4.

\subsection{Questão 1 — PID vs. PPID}

\textit{Qual a diferença entre PID e PPID? Ilustre com um exemplo observado na Atividade 1.3.}

\textbf{Resposta:} O PID é o identificador único do próprio processo, atribuído pelo kernel no momento da criação; o PPID é o identificador do processo \textbf{pai}, ou seja, daquele que o criou. O primeiro responde ``quem sou eu'', o segundo responde ``quem me criou''.

O exemplo observado na Atividade 1.3 é a cadeia \texttt{explorer.exe} (PID 4328) $\rightarrow$ \texttt{cmd.exe} (PID 10552) $\rightarrow$ \texttt{python.exe} (PID 11204). O \texttt{python.exe} tem PID 11204 e PPID 10552, valor que corresponde exatamente ao PID do \texttt{cmd.exe} de onde o script foi lançado. É essa correspondência entre o PPID de um processo e o PID de outro que encadeia a hierarquia em árvore do sistema.

\subsection{Questão 2 — Múltiplas instâncias e compartilhamento de memória}

\textit{Ao abrir três instâncias do mesmo programa, quantos PIDs você observou? Eles compartilham memória entre si?}

\textbf{Resposta:} Foram observados \textbf{três PIDs distintos} — 8452, 9128 e 9864 (Tabela \ref{tab:pids}) —, um para cada janela do Bloco de Notas, embora as três executem o mesmo arquivo em disco (\texttt{C:\textbackslash Windows\textbackslash System32\textbackslash notepad.exe}).

\textbf{Não}, esses processos não compartilham memória entre si. Cada um recebeu do sistema operacional seu próprio espaço de endereçamento virtual, isolado dos demais: o texto digitado em uma janela não aparece nas outras, e uma falha em uma instância não derruba as demais. Esse isolamento é justamente a garantia de proteção que caracteriza o processo como unidade.

\subsection{Questão 3 — Número de threads observado}

\textit{Quantas threads o processo analisado na Atividade 1.2 possuía? Isso foi surpreendente?}

\textbf{Resposta:} O processo Python analisado possuía \textbf{apenas 1 thread} (Tabela \ref{tab:threads}), a thread principal, com o identificador 12048.

\textbf{Sim, o resultado foi surpreendente} em um primeiro momento: esperava-se que o interpretador Python mantivesse threads auxiliares próprias. A investigação posterior mostrou a razão: por causa do GIL (\textit{Global Interpreter Lock}), o CPython executa \textit{bytecode} em uma única thread por vez e só cria threads adicionais quando o programa as solicita explicitamente pelo módulo \texttt{threading}. O contraste com o Google Chrome, que apresentou 38 threads no mesmo instante, evidencia como a contagem depende inteiramente da arquitetura da aplicação, e não do sistema operacional.

\subsection{Questão 4 — Concordância entre psutil e a ferramenta nativa}

\textit{Os resultados obtidos pelo script Python (psutil) bateram com os resultados da ferramenta nativa (ps/Gerenciador de Tarefas) em todas as atividades? Alguma diferença encontrada?}

\textbf{Resposta:} \textbf{Sim, em todas as quatro atividades}, conforme consolidado na Tabela \ref{tab:comparacao}. Nas Atividades 1.1, 1.2 e 1.3 os valores foram idênticos dígito a dígito: os três PIDs do Bloco de Notas listados pelo script coincidiram com os exibidos na aba Detalhes do Gerenciador de Tarefas; as contagens de threads retornadas pelo \texttt{psutil} reproduziram as do \texttt{Get-Process}; e a cadeia de ancestrais reconstruída pelo script repetiu os pares PID/PPID lidos no \texttt{Get-CimInstance Win32\_Process}.

\textbf{Uma única diferença foi encontrada}, na Atividade 1.4, e ela é de \textbf{unidade de medida, não de conteúdo}: o \texttt{monitor\_cpu.py} indicava $\sim$100\%, enquanto o Gerenciador de Tarefas mostrava $\sim$8\% para o mesmo processo no mesmo instante. A explicação é a base de normalização — o \texttt{cpu\_percent()} do \texttt{psutil} toma um núcleo lógico como referência, ao passo que o Gerenciador expressa o consumo em relação à capacidade total da máquina. Como o processador possui 12 núcleos lógicos, $100\% \div 12 \approx 8.3\%$, e as duas leituras descrevem exatamente o mesmo fenômeno. A lição prática é que comparar ferramentas exige verificar antes qual é a base de cada percentual.

\subsection{Questão 5 — Ferramentas utilizadas e obstáculos}

\textit{Qual ferramenta você usou (linha de comando nativa, gráfica, ou Python) e por quê? Descreva um obstáculo prático encontrado na instalação/uso.}

\textbf{Resposta:} Foram usadas as \textbf{três} categorias, de forma complementar: a ferramenta \textbf{gráfica} (Gerenciador de Tarefas) para a leitura rápida de PIDs e para a visualização do gráfico de CPU; a \textbf{linha de comando nativa} (PowerShell, com \texttt{Get-Process} e \texttt{Get-CimInstance Win32\_Process}) para obter listas de threads e pares PID/PPID em formato tabular; e o \textbf{Python com psutil} como referência principal, por duas razões: o mesmo código roda sem alteração em Windows e Linux, e a saída fica registrada em texto, o que facilita documentar os resultados neste relatório.

\textbf{Obstáculo prático encontrado:} a instalação inicial do \texttt{psutil}. O comando \texttt{pip} não era reconhecido no terminal porque o Python havia sido instalado sem marcar a opção \textit{Add python.exe to PATH}; foi necessário corrigir a variável de ambiente PATH antes que \texttt{pip install psutil} funcionasse. Um segundo ponto de atenção foi a Atividade 1.4, que exige \textbf{dois terminais abertos simultaneamente} — no primeiro o gerador de carga, no segundo o monitor —, já que executar os dois em sequência no mesmo terminal não permite observar o consumo enquanto ele ocorre.

\subsection{Questão 6 — Diferença entre processo e thread}

\textit{Com base na experiência prática, escreva com suas próprias palavras a diferença entre processo e thread.}

\textbf{Resposta:} Com base no que foi observado nas quatro atividades, um \textbf{processo} é uma \textbf{unidade de isolamento}: é o programa efetivamente carregado na memória, com PID próprio, espaço de endereçamento próprio e recursos próprios. Foi o que se viu na Atividade 1.1, em que três janelas do mesmo executável se comportaram como três mundos separados, sem enxergar a memória umas das outras.

Uma \textbf{thread} é uma \textbf{unidade de execução dentro de um processo}: é um fluxo de instruções que o sistema operacional agenda para rodar. As threads de um mesmo processo compartilham a memória desse processo, mas cada uma mantém seus próprios registradores e sua própria pilha. Foi o que se viu na Atividade 1.2, em que um único Chrome, com um único PID, abrigava 38 fluxos de execução simultâneos.

Em resumo: \textbf{processos separam, threads compartilham}. Usam-se múltiplos processos quando se quer isolamento e segurança (o travamento de um não afeta os outros); usam-se múltiplas threads quando se quer eficiência e troca barata de dados dentro de uma mesma aplicação. O preço de cada escolha é o simétrico do seu benefício: comunicação entre processos é mais custosa, e compartilhamento entre threads exige sincronização cuidadosa.

\subsection{Considerações finais}

Os experimentos comprovaram que ferramentas de observação do SO — tanto nativas quanto via APIs como o \texttt{psutil} — permitem validar e aprofundar o entendimento teórico de conceitos de sistemas operacionais. A discrepância entre a expectativa teórica (múltiplas threads no processo Python) e a prática (apenas 1 thread observada) levou a uma investigação mais profunda sobre o GIL, demonstrando que a observação empírica é essencial para complementar o conhecimento teórico.

\section{REFERÊNCIAS BIBLIOGRÁFICAS}

\noindent [1] Tanenbaum, A. S., \textit{"Modern Operating Systems"}, 4ª edição, Pearson, 2015.

\noindent [2] Silberschatz, A., Galvin, P. B., Gagne, G., \textit{"Operating System Concepts"}, 10ª edição, Wiley, 2018.

\noindent [3] Documentação oficial psutil, \textit{"psutil documentation"}, disponível em: https://psutil.readthedocs.io/, acesso em agosto de 2026.

\noindent [4] Microsoft Docs, \textit{"About Processes and Threads"}, disponível em: https://docs.microsoft.com/windows/win32/procthread/, acesso em agosto de 2026.

\end{document}
